\documentclass[11pt]{article}

% ---------- packages ----------
\usepackage[utf8]{inputenc}
\usepackage[T1]{fontenc}
\usepackage{amsmath, amssymb, amsthm}
\usepackage[dvipsnames, svgnames]{xcolor} 
\usepackage{graphicx}
\usepackage{hyperref}
\usepackage{booktabs}
\usepackage{microtype}
\usepackage{geometry}
\geometry{a4paper, margin=2.5cm}

% ---------- hyperref config ----------
\hypersetup{
  colorlinks = true,
  linkcolor  = blue!60!black,
  citecolor  = teal,
  urlcolor   = teal,
  pdftitle   = {Entropic Semantic Diffusion: Spectral Search and Internal Time from Entropy},
  pdfauthor  = {Lorenzo Moriondo},
}

% ---------- theorem environments ----------
\newtheorem{definition}{Definition}
\newtheorem{proposition}{Proposition}
\newtheorem{conjecture}{Conjecture}
\newtheorem{remark}{Remark}

% ---------- custom commands ----------
\newcommand{\ESDM}{\textsc{esdm}}
\newcommand{\VDT}{\textsc{vdt}}
\newcommand{\SDM}{\textsc{sdm}}
\newcommand{\SPIN}{\textsc{spin}}
\newcommand{\Ltau}{\Lambda(\tau)}
\newcommand{\taun}{\tau_n}
\newcommand{\alphabar}{\bar{\alpha}(\tau)}
\newcommand{\Lf}{L_F}

% ============================================================
\begin{document}
% ============================================================

\title{%
  \textbf{Entropic Semantic Diffusion Model:\\Spectral Search and Internal Time from Entropy}%
}

\author{%
  Lorenzo Moriondo\\[2pt]
\small Independent Researcher --- tuned.org.uk\\
\small ORCID: 0000-0002-8804-2963\\
  \texttt{tunedconsulting@gmail.com} \\
  \href{https://github.com/tuned-org-uk/enthropic-semantic-diffusion}{%
    github.com/tuned-org-uk/enthropic-semantic-diffusion}}


\date{June 2026 \\ \small{Draft -- comments welcome}}

\maketitle

% ============================================================
\begin{abstract}
% ============================================================

Diffusion models and variational sequence architectures parametrise their
dynamics by an \emph{external} time variable -- a fixed integer step counter or
a predetermined noise schedule -- that is semantically meaningless and
thermodynamically inert.
This is a semantic instance of the \emph{Wheeler--DeWitt problem of time}: a
closed system described by a time-independent Hamiltonian admits no intrinsic
mechanism to order its own evolution.
We draw on a recent cold-atom experiment by Barontini (2026), in which a
Bose--Einstein condensate partitioned into an observed (\emph{bright}) and an
unobserved (\emph{dark}) sector is shown to self-sequence its own dynamics via
\emph{entropic time} -- a monotonically increasing quantity constructed entirely
from internal entropy exchange between the two sectors -- and proposes an
effective Schr\"odinger equation parametrised by this internal time that
quantitatively reproduces the measured dynamics.

The \emph{Entropic Semantic Diffusion Model} (\ESDM) applies this lesson to a
unified semantic architecture combining the Vibrational Deductive Transformer
(\VDT), the Spectral Diffusion Model (\SDM), and the ArrowSpace dual-space
hypothesis.
The corpus embedding space is partitioned into two sectors: the \emph{bright
sector} of active spectral modes $\omega_k > \theta$ (observable, governs
reconstruction) and the \emph{dark sector} of latent feature-space wiring
$L_F$ and inactive modes (unobserved, governs dynamics).
Entropic time $\taun = \sum_i |\Delta\,\text{kl\_tau}_i|$ accumulates from the
absolute change in the tau-mode KL divergence of the \VDT ELBO; it is
monotonically increasing by construction and stops advancing when no entropy is
exchanged between sectors -- the semantic analogue of Barontini's heat death.
The fixed cosine diffusion schedule is replaced by $\alphabar =
\cos^2\!\bigl(\frac{\pi}{2}\cdot\frac{\tau}{\tau_{\max}}\bigr)$, so that
semantic structure dissolves at precisely the moments the spectral geometry is
actively reorganising.
The \VDT wave recurrence is gated by an \emph{entropic pump}
$\Ltau = |\partial_\phi\,\text{kl\_tau}|^{-1}$ that amplifies attention forcing
during low-entropy-flow regimes and suppresses it during rapid mode transitions,
in direct correspondence with Barontini's Eq.~(6).
An entropy conservation term $\beta\,|S_{\text{bright}}(\tau) +
S_{\text{dark}}(\tau) - S_0|^2$ is added to the ELBO, enforcing the global
constancy of semantic entropy observed empirically in the BEC miniuniverse.

The resulting architecture has no external epoch counter, no preset number of
diffusion steps, and no fixed noise schedule: its temporal structure emerges
entirely from the internal thermodynamics of the dual-space semantic system.
Training phases (big bang, expansion, maximum extension, big crunch, heat
death) map onto identifiable regimes of the spectral diagnostics, and a
\texttt{kl\_S} barrier-height sweep directly reproduces the topological
transition structure of Barontini's Fig.~2.
We argue that this construction formalises the ArrowSpace dual-space
hypothesis -- every vector space is simultaneously a geometric item-space and
a semantic feature-space -- as a thermodynamic duality, and that optimisation
in the \ESDM is equivalently a \emph{recursive semantic search} over the
spectral geometry of the corpus.

\end{abstract}

\vspace{1em}
\noindent\textbf{Keywords:} entropic time, spectral diffusion, vibrational transformer,
Wheeler--DeWitt problem, ArrowSpace, graph Laplacian, semantic search,
bose--einstein condensate, variational inference.

\newpage

% ============================================================
\tableofcontents
% ============================================================

\newpage

% ============================================================
\section{Introduction}
\label{sec:introduction}
% ============================================================
% TODO: expand
% - The semantic gap and the dual-space hypothesis (ArrowSpace / SPIN)
% - External time as a design limitation in VDT and SDM
% - The Wheeler-DeWitt problem as the unifying framing
% - Barontini as experimental existence proof
% - Paper contributions and structure

\section{Background}
\label{sec:background}

  \subsection{The ArrowSpace Dual-Space Hypothesis}
  \label{sec:arrowspace}
  % TODO:
  % - Item-space vs. feature-space Laplacian L_F
  % - Graph Wiring (GW) and epiplexity
  % - Fiedler vector as semantic scale factor
  % - SPIN: spectral augmentation of retrieval

  \subsection{Vibrational Deductive Transformer (VDT)}
  \label{sec:vdt}
  % TODO:
  % - Wave recurrence, VibrationalStateBlock
  % - SignedDensityMatrix, rho_+, rho_-
  % - ELBO: reconstruction + kl_z + kl_S + kl_tau
  % - External time flaw

  \subsection{Spectral Diffusion Model (SDM)}
  \label{sec:sdm}
  % TODO:
  % - LGMRF prior, per-mode variance
  % - cosine_alpha_bar(t, T=200) and its limitations
  % - Connection to VDT latent structure

  \subsection{Barontini's Experiment: Internal Time from Entropy}
  \label{sec:barontini}
  % TODO:
  % - BEC setup, bright / dark partition, optical barrier
  % - Entropic time definition (Eq. 3)
  % - Entropy conservation law
  % - Entropic Schrödinger equation (Eq. 6): phi, Lambda(tau)
  % - Three key experimental results

\section{The Wheeler--DeWitt Problem in Semantic Architectures}
\label{sec:wdw}
% TODO:
% - Formal statement: VDT + SDM as a closed system with external t
% - Why external time is a structural flaw, not just a design choice
% - The two-sector partition as the solution structure

\section{Entropic Time for Semantics}
\label{sec:entropic_time}

  \subsection{Definition and Monotonicity}
  \label{sec:tau_definition}
  % tau_n = sum |Delta kl_tau_i|, monotonicity proof

  \subsection{Clock Field and Non-Monotonicity Resolution}
  \label{sec:clock_field}
  % log_lam_s as phi, absolute-value integral fix

  \subsection{Heat Death and Termination}
  \label{sec:heat_death}
  % tau_n stalling as intrinsic stopping criterion

\section{Entropic Semantic Diffusion Model (ESDM)}
\label{sec:esdm}

  \subsection{Two-Sector Partition}
  \label{sec:partition}
  % Bright: N_active modes; Dark: L_F, inactive modes
  % Full correspondence table (Barontini vs. ESDM)

  \subsection{The Entropic Pump and Gated Wave Recurrence}
  \label{sec:wave_recurrence}
  % Lambda(tau) = |d_phi kl_tau|^{-1}
  % Modified VibrationalStateBlock equation
  % Interpretation: exploration vs. free dynamics

  \subsection{Entropic Diffusion Schedule}
  \label{sec:schedule}
  % alpha_bar(tau), per-mode schedule alpha_bar_k(tau_k)
  % Comparison to cosine_alpha_bar(t, T)

  \subsection{The ESDM ELBO}
  \label{sec:elbo}
  % Full five-term objective
  % Entropy conservation term: beta calibration

  \subsection{Signed Interference as Confidence}
  \label{sec:confidence}
  % rho_sign = rho_- / (rho_+ + eps)
  % Low-confidence detection without external validation

  \subsection{Architecture Summary}
  \label{sec:architecture}
  % Four modules: EntropicClock, VibrationalStateBlock,
  %               SpectralLoadingDecoder, EntropicDiffusionDecoder
  % Data flow diagram

\section{Training Dynamics and Diagnostics}
\label{sec:training}

  \subsection{Per-Batch Health Monitoring}
  \label{sec:monitoring}
  % kl_bright, kl_dark, tau_n, Lambda_tau, conserved, rho_sign
  % Health check table (healthy / warning / critical)

  \subsection{Phase Identification}
  \label{sec:phases}
  % Big bang, expansion, maximum extension, big crunch, heat death
  % Barontini correspondence

\section{Experimental Protocol}
\label{sec:experiments}

  \subsection{Barrier Height Sweep (kl\_S analogue of $V$)}
  \label{sec:barrier_sweep}
  % kl_S sweep: tau_n vs. training-step curves
  % Predicted regimes: cyclic / clean expansion / heat death
  % Direct analogue of Barontini Fig. 2

  \subsection{Entropy Conservation Verification}
  \label{sec:conservation_exp}
  % S_bright(tau) + S_dark(tau) = S_0 throughout training

  \subsection{Loschmidt Echo: Reversibility Test}
  \label{sec:loschmidt}
  % Run tau backwards; measure embedding recovery
  % Arrow-of-time test in semantic space

\section{Connection to Thermal Time and Recursive Semantic Search}
\label{sec:thermal_time}

  \subsection{Connes--Rovelli Thermal Time Hypothesis}
  \label{sec:connes_rovelli}
  % Operational vs. algebraic construction of time

  \subsection{ArrowSpace Dual-Space as Thermodynamic Duality}
  \label{sec:duality}
  % Item-space = bright sector; feature-space = dark sector
  % Spectral Complementarity Conjecture

  \subsection{Optimisation as Recursive Semantic Search}
  \label{sec:recursive_search}
  % L_F -> dark entropy -> tau -> schedule -> loss -> L_F loop
  % SPIN as the retrieval-time projection of this loop

\section{Open Questions}
\label{sec:open_questions}
% 1. Feshbach-Villars decomposition in the discrete case
% 2. Multiple clock fields: log_lam_s vs. Fiedler eigenvalue
% 3. Quantum bounce vs. singularity: sign flip of Lambda
% 4. Loschmidt echo as formal reversibility test
% 5. beta calibration for the conservation term

\section{Conclusion}
\label{sec:conclusion}
% TODO

% ============================================================
\newpage
\begin{thebibliography}{9}

\bibitem{Barontini2026}
G.~Barontini,
\newblock {Testing the problem of time with cold atoms},
\newblock {\em Physical Review Research}, 8:L022047, 2026.
\newblock doi: \href{https://doi.org/10.1103/1h9j-df4k}{10.1103/1h9j-df4k}.

\bibitem{MoriondoAzizi2026}
L.~Moriondo and I.~Azizi,
\newblock {From Embedding Geometry to Spectral Search: Energy Dispersion Networks for Vector Retrieval},
\newblock arXiv:2606.21535 [cs.IR], 2026.
\newblock \href{https://arxiv.org/abs/2606.21535}{https://arxiv.org/abs/2606.21535}.

\bibitem{Moriondo2025JOSS}
L.~Moriondo,
\newblock {ArrowSpace: introducing Spectral Indexing for vector search},
\newblock {\em Journal of Open Source Software}, 10(113):9002, 2025.
\newblock doi: \href{https://doi.org/10.21105/joss.09002}{10.21105/joss.09002}.

\end{thebibliography}
% ============================================================
% refs.bib entries to include:
%  Barontini2026, ConnesRovelli1994, PageWootters1983,
%  KieferSandhofer2022, DeWitt1967,
%  Moriondo2025VDT, Moriondo2025SDM, Moriondo2025QV,
%  Moriondo2024ArrowSpace, Moriondo2026SPIN

\end{document}
